\documentclass[11pt,a4paper]{article}
\usepackage[margin=1in]{geometry}
\usepackage{booktabs}
\usepackage{hyperref}
\usepackage{xurl}
\usepackage{amsmath}

\title{Step5a-HGAT Paper Preparation (Detailed, EN)}
\author{}
\date{\today}

\begin{document}
\maketitle

\section{Scope and Mandatory Protocol}
This draft follows a single mandatory protocol:
\begin{itemize}
  \item target split mode: \texttt{table\_group};
  \item split ratio: \textbf{train:val = 1:5};
  \item \textbf{no test split} (test=0).
\end{itemize}

Dataset:
\begin{center}
\path{output/dataset_table_group_train_val_no_test.pkl}
\end{center}
Target-domain counts: train=1274, val=6270, test=0.

\section{Algorithm Core}
We use Topology-Conditioned Dual Residual (TCDR) on top of a shared backbone:
\begin{align}
\mathbf{h} &= f_{\theta}([\mathbf{x};\mathbf{z}]), \\
\hat{y}_{base} &= \mathbf{w}^{\top}\mathbf{h} + b, \\
\Delta_{src} &= g_{\phi_{src}}([\mathbf{h};\mathbf{e}_t]), \\
\Delta_{tgt} &= g_{\phi_{tgt}}([\mathbf{h};\mathbf{e}_t]).
\end{align}
Final prediction is base plus domain-specific residual.

Training objective:
\begin{equation}
\mathcal{L} = \mathcal{L}_{tgt} + \lambda(e)\mathcal{L}_{src}.
\end{equation}

\section{No-Test Result Summary}
The following tables are generated from:
\begin{itemize}
  \item \path{paper_materials/tables/paper_notest_core_summary.csv}
  \item \path{paper_materials/tables/paper_t32_onefactor_summary.csv}
\end{itemize}

\subsection{Core Multi-Seed Comparison}
\begin{table}[htbp]
\centering
\small
\begin{tabular}{lcccc}
\toprule
Method & NSeed & Train $R^2$ (mean+/-std) & Val $R^2$ (mean+/-std) & Val Loss (mean+/-std) \\
\midrule
Baseline(shared linear) & 5 & 0.9024+/-0.0086 & 0.8726+/-0.0067 & 0.121635+/-0.006428 \\
TCDR(16,96) & 5 & 0.9586+/-0.0085 & 0.9065+/-0.0168 & 0.089249+/-0.016086 \\
TCDR(32,64) & 5 & \textbf{0.9582+/-0.0216} & \textbf{0.9234+/-0.0175} & \textbf{0.073105+/-0.016755} \\
\bottomrule
\end{tabular}
\caption{Core architecture comparison under the 1:5 no-test protocol.}
\end{table}

\subsection{One-Factor Ablation (Main Backbone: T32x64)}
\begin{table}[htbp]
\centering
\small
\begin{tabular}{lcccc}
\toprule
Method & NSeed & Train $R^2$ (mean+/-std) & Val $R^2$ (mean+/-std) & Val Loss (mean+/-std) \\
\midrule
T32x64 + mode\_base(enrich\_on) & 5 & 0.9682+/-0.0087 & \textbf{0.9298+/-0.0073} & \textbf{0.066996+/-0.006930} \\
T32x64 + enrich\_off(auto) & 5 & 0.9682+/-0.0087 & \textbf{0.9298+/-0.0073} & \textbf{0.066996+/-0.006930} \\
T32x64 + parasitic\_replace & 5 & 0.9694+/-0.0080 & 0.9294+/-0.0084 & 0.067451+/-0.008056 \\
T32x64 + dual\_readout(mean) & 5 & 0.9582+/-0.0216 & 0.9234+/-0.0175 & 0.073105+/-0.016755 \\
T32x64 + mode\_base+topology\_off & 5 & 0.9006+/-0.0068 & 0.8720+/-0.0034 & 0.122180+/-0.003232 \\
\bottomrule
\end{tabular}
\caption{One-factor ablation on T32x64 under the same protocol.}
\end{table}

\section{Current Conclusion for Writing}
\begin{itemize}
  \item TCDR is the major contributor compared with baseline.
  \item Under the current no-test setting, TCDR(32,64) is the strongest candidate.
  \item Under fixed T32x64, mode\_base(enrich\_on) and enrich\_off(auto) are strongest and numerically identical.
  \item Topology-off causes a large drop (Val $R^2$: 0.9298 $\rightarrow$ 0.8720), supporting topology-conditioned residual as the key contribution.
  \item T16x96 one-factor ablations are kept as supplementary comparison (not the main table).
\end{itemize}

\section{Reproducibility}
Statistics are reproducible by:
\begin{center}
\path{ICCAD2026_Changxin/paper_materials/tables/build_step5a_notest_tables.ps1}
\end{center}
and the T32 one-factor table is rebuilt by:
\begin{center}
\path{ICCAD2026_Changxin/paper_materials/tables/build_step5a_t32_onefactor_tables.ps1}
\end{center}

\end{document}
